% =====================================================================
%  Label-flipping attack surface: the reachable deviation set
%  Self-contained section, unified with the existing framework.
%  No new \newcommand: only existing macros are used.
%
%  STRUCTURE
%    Part 0  Setup: surface, shifts, budget, kernels, federated GGUs,
%            the three admissible sets, one-shot poisoning
%    Part A  Mean aggregation: closed-form reachable set, geometry,
%            feasibility, AND the attacker's problem
%    Part B  Robust aggregation: coordinate-wise selection framework,
%            exact decomposition of Agg - Mean, surrogate optimality,
%            stealth and selection
%    Part C  Validation TODO
% =====================================================================

% ---------------------------------------------------------------------
% NOTATION RECONCILIATION
%   1. \delta is the trigger; the gradient gap is \norm{v_k}
%      (\delta(\thk) is retired from lem:bdagger-kappa and
%       prop:feasibility-kappa, where it is a local shorthand).
%   2. \mathcal{C}_k is the event \{\Npk\le f\}; the flip polytope is
%      \mathcal{F}(\thk).
%   3. The budget fraction is \beta; the label count is N_{flip}
%      (b_k is the induced bias, \mathcal{B}_k the bias map).
%   4. The selection size of an aggregation rule is \ell.
%   5. Class indices are (y,z); generic vectors of \R^d are g, e, p.
%   6. \varpi_k is the RANK ratio (\pi_y is a class probability).
%   7. THE SHIFT VECTORS CARRY NO \pi_y FACTOR: see def:shifts and
%      rem:no-pi. This corrects an earlier draft in which
%      \bar G_{y,z} was defined as \pi_y(g_{y,z}-g_{y,y}), which is
%      inconsistent with u_{y,z} being a joint proportion.
%   Adopted unchanged: b_k, \mathcal{E}_k, v_k, a_k.
% ---------------------------------------------------------------------


\subsection{Setting}
\label{sec:setting}

The defender solves
\begin{equation}
\min_{\theta\in\Th}\ \Lc(\theta):=\mathbb{E}_{(x,y)\sim\Dc}\big[\ell_c(f_\theta(x),y)\big],
\end{equation}
where $\Dc$ is a distribution over $\mathcal X\times\mathcal Y$, $\Th\subset\R^d$ and $\ell_c:\mathcal X\times\mathcal Y\to\R$. The attacker solves
\begin{equation}
\min_{\theta\in\Th}\ \Lp(\theta):=\mathbb{E}_{(x,y)\sim\Dp}\big[\ell_c(f_\theta(x),y)\big],
\end{equation}
with $\Dp$ chosen according to the attacker's goal, as in \cref{ex:backdoor}.

\begin{example}[Backdoor attack]
\label{ex:backdoor}
Let $T:\mathcal X\to\mathcal X$ be a measurable \emph{trigger} map and $y_{\mathrm{source}}\neq y_{\mathrm{target}}\in\mathcal Y$ a source/target pair. The associated \emph{backdoor operator} is
\[
\Phi:\mathcal X\times\mathcal Y\to\mathcal X\times\mathcal Y,
\qquad
\Phi(x,y)=
\begin{cases}
\big(T(x),y_{\mathrm{target}}\big), & y=y_{\mathrm{source}},\\
(x,y), & \text{otherwise.}
\end{cases}
\]
Assuming $p_s:=\mathbb{P}_{\Dc}(Y=y_{\mathrm{source}})>0$, write $\mathcal D_s:=\Dc(\cdot\mid Y=y_{\mathrm{source}})$. Since $\mathcal D_s$ is concentrated on $\mathcal X\times\{y_{\mathrm{source}}\}$, the pushforward $\Phi_\#\mathcal D_s$ is the law of $(T(X),y_{\mathrm{target}})$ with $(X,Y)\sim\mathcal D_s$. The attacker targets the mixture
\begin{equation}\label{eq:bd-Dp}
\Dp=(1-\lambda)\,\Dc+\lambda\,\Phi_\#\mathcal D_s,\qquad\lambda\in(0,1),
\end{equation}
so that
\begin{equation}\label{eq:bd-Lp-split}
\Lp(\theta)=(1-\lambda)\,\Lc(\theta)+\lambda\,\mathbb{E}_{(X,Y)\sim\mathcal D_s}\Big[\ell_c\big(f_\theta(T(X)),y_{\mathrm{target}}\big)\Big].
\end{equation}
\end{example}


% =====================================================================
\subsection*{Part 0 --- Setup}
% =====================================================================

\subsection{The label-flipping attack surface}
\label{sec:flip-surface}

We consider an attacker able to invert a constrained budget of labels but unable to alter inputs.

\begin{definition}[Label-flipping attack surface]
\label{def:flip-surface}
\begin{equation}
\attsuf{\thk}
=\Big\{g\in\R^d\;\Big|\;\exists\,(x,y)\in\operatorname{supp}\Dc,\ \exists\,c\in\mathcal{Y},\
g=\nabla_{\thk}\ell_c\big(f_{\thk}(x),c\big)\Big\}.
\end{equation}
\end{definition}

\begin{definition}[Class-conditional gradients, shifts, shift radius]
\label{def:shifts}
For $y\neq z$ in $\mathcal Y$, with $C:=|\mathcal{Y}|$ and $\pi_y:=\mathbb{P}_{\Dc}(Y=y)$,
\begin{equation}
g_{y,z}(\thk):=\nabla_\theta\,\mathbb{E}_{X\mid Y=y}\big[\ell_c(f_{\thk}(X),z)\big],
\qquad
\bar G_{y,z}(\thk):=g_{y,z}(\thk)-g_{y,y}(\thk),
\label{eq:shift}
\end{equation}
and $\bar G(\thk)\in\R^{d\times C(C-1)}$ collects the \emph{shift vectors} $\bar G_{y,z}$ as
columns. We write
\begin{equation}
\varsigma_k:=\max_{z\neq y}\norm{\bar G_{y,z}(\thk)}
\label{eq:varsigma}
\end{equation}
for the \emph{shift radius} at round $k$. Note $\gradc(\thk)=\sum_y\pi_yg_{y,y}(\thk)$.
\end{definition}

\begin{definition}[Flip polytope]
\label{def:Fpoly}
$\mathcal{F}(\thk):=\mathrm{conv}\big(\{0\}\cup\{\bar G_{y,z}(\thk)\}_{z\neq y}\big)\subset\R^d$.
\end{definition}

$\mathcal{F}(\thk)$ is the convex hull of at most $C(C-1)+1$ points, hence a compact convex
polytope independent of $\beta$ and of the trigger. Geometrically $\bar G_{y,z}(\thk)$ is the
deviation produced by relabeling a \emph{unit mass} of class-$y$ samples into $z$;
$\mathcal{F}(\thk)$ collects the deviations obtainable by mixing such relabelings, with $0$
available as ``do nothing''.


\subsection{Write access: budget, allocation, corruption rates}
\label{sec:budget}

\begin{definition}[Budget, allocation, corruption rates]
\label{def:budget}
Let $\mathcal{M}\subseteq[n_b]$, $|\mathcal{M}|=n_p$, be the malicious units. Shards are of equal size $n/n_b$, and $N_{\mathrm{flip}}\in\N$ is the number of labels the adversary may \emph{alter} (a sample selected for poisoning but relabeled to its own class does not count).
\begin{itemize}[noitemsep,leftmargin=*]
\item \textbf{Global budget.} $\beta:=N_{\mathrm{flip}}/n\in(0,1)$, the fraction of \emph{all} training labels altered. This is the primitive constraint.
\item \textbf{Allocation.} Unit $i$ receives $N^i_{\mathrm{flip}}$ alterations, $N^i_{\mathrm{flip}}=0$ for $i\notin\mathcal{M}$ and $\sum_{i\in\mathcal{M}}N^i_{\mathrm{flip}}=N_{\mathrm{flip}}$.
\item \textbf{Local corruption rate.} $\beta^i:=N^i_{\mathrm{flip}}\big/(n/n_b)$, the fraction of unit $i$'s \emph{own} shard altered; admissibility requires $\beta^i\le1$.
\item \textbf{Average corruption rate.} $\bar\beta:=\tfrac1{n_b}\sum_i\beta^i$.
\end{itemize}
\end{definition}

\begin{lemma}[The average corruption rate is the global budget]
\label{lem:beta-bar}
For every admissible allocation $\bar\beta=\beta$, and $\beta^i\le1$ for all $i$ forces
\begin{equation}
\beta\ \le\ \gamma:=\frac{n_p}{n_b}.
\label{eq:beta-feasible}
\end{equation}
\end{lemma}
\begin{proof}
$\bar\beta=\tfrac1{n_b}\sum_{i\in\mathcal{M}}N^i_{\mathrm{flip}}\tfrac{n_b}{n}=N_{\mathrm{flip}}/n=\beta$; and $N_{\mathrm{flip}}\le n_p\,n/n_b$ gives \cref{eq:beta-feasible}.
\end{proof}

We write $\gamma=n_p/n_b$ throughout. Uniform allocation gives $\beta^i=\beta/\gamma$ on
$\mathcal{M}$ and $0$ elsewhere, and \cref{eq:beta-feasible} is exactly the condition that this
be $\le1$. Since $\beta^i=0$ off $\mathcal{M}$, at most $n_p$ units hold altered labels;
\cref{def:ggu-fed} fixes the convention by which $\Npk$ counts them.


\subsection{Heterogeneous flipping kernels}
\label{sec:kernels}

The modelled attack produces a poisoned dataset \emph{before} the victim's training begins, so
the attack policy of each worker must be parameterized independently of the current round
$k\in\N$. We therefore parameterize the attack policy of each malicious worker $i$ by a single
vector $u^i$, fixed once and for all.

\begin{definition}[Flipping kernel, selection rates, flip masses]
\label{def:kernel}
Unit $i$ applies a \emph{flipping kernel} $w^i$, with $w^i_{y,z}\ge0$ the probability that a class-$y$ sample \emph{selected for poisoning} is presented with label $z$ and $\sum_zw^i_{y,z}=1$; and \emph{class-dependent selection rates} $\beta^i_y\in[0,1]$, the fraction of unit $i$'s class-$y$ samples selected. The induced \emph{flip masses} are
\begin{equation}
u^i_{y,z}:=\pi_y\,\beta^i_y\,w^i_{y,z}\quad(z\neq y),
\qquad u^i:=(u^i_{y,z})_{z\neq y}\in\R^{C(C-1)},
\end{equation}
with $u^i=0$ for $i\notin\mathcal{M}$. The associated policy $\ga^i(\thk)$ is the law of $\nabla_{\thk}\ell_c(f_{\thk}(X),Z)$ under this relabeling.
\end{definition}

\begin{remark}[Reading $u^i$]
\label{rem:reading-u}
The vector $u^i$ reads as follows.
\begin{itemize}
\item \emph{Coordinates.} $u^i_{y,z}$ is a \emph{joint} proportion: among the samples held by unit $i$, the fraction truly of class $y$ \emph{and} carrying label $z$ after poisoning. Nonnegative as a product of a probability, a rate and a kernel entry.
\item \emph{Total mass.} $\|u^i\|_1=\sum_y\pi_y\beta^i_y(1-w^i_{y,y})=\beta^i$, consistently with \cref{def:budget}: only \emph{altered} labels are counted.
\item \emph{Per-class mass.} $\sum_{z\neq y}u^i_{y,z}\le\pi_y$: one cannot relabel more class-$y$ samples than the shard contains. This \emph{capacity} constraint is the only structural restriction.
\item \emph{Wasted budget.} A diagonal entry $w^i_{y,y}>0$ means a sample was selected but kept its label; it consumes selection without producing a gradient shift, whence only off-diagonal indices appear.
\end{itemize}
\end{remark}

\begin{lemma}[Expected message]
\label{lem:message}
Under \cref{def:kernel}, for every unit $i$ and every $k$,
\[
\Ethk{g_i(\thk)}=\gradc(\thk)+\bar G(\thk)u^i.
\]
\end{lemma}
\begin{proof}
The expected gradient emitted by unit $i$ is
$\sum_y\pi_y[(1-\beta^i_y)g_{y,y}+\beta^i_y\sum_zw^i_{y,z}g_{y,z}]$. Subtract
$\gradc(\thk)=\sum_y\pi_yg_{y,y}$ and use $\sum_zw^i_{y,z}=1$ to obtain
$\sum_y\sum_{z\neq y}\pi_y\beta^i_yw^i_{y,z}(g_{y,z}-g_{y,y})
=\sum_y\sum_{z\neq y}u^i_{y,z}\,\bar G_{y,z}(\thk)=\bar G(\thk)u^i$,
the middle equality being \cref{def:kernel} together with \cref{eq:shift}.
\end{proof}

\begin{remark}[Why the shift vectors carry no $\pi_y$]
\label{rem:no-pi}
The proof of \cref{lem:message} is what fixes the convention. Since $u^i_{y,z}$ already contains
the factor $\pi_y$ --- it is a proportion of the \emph{whole} shard, not of class $y$ alone
(\cref{rem:reading-u}) --- the columns of $\bar G$ must be the bare differences
$g_{y,z}-g_{y,y}$. Defining $\bar G_{y,z}:=\pi_y(g_{y,z}-g_{y,y})$ instead, as an earlier draft
did, double-counts $\pi_y$ and understates every induced deviation by that factor: on a balanced
$C$-class problem, by a factor $C$. The capacity constraint $\sum_{z\neq y}u_{y,z}\le\pi_y$ and
the identity $\|u^i\|_1=\beta^i$ are the two invariants that must be preserved; both are
consistent with \cref{eq:shift} and with it alone. Note the practical consequences: $\varsigma_k$,
$\varrho_k$ and the certificates of \cref{prop:certificates-mean} all scale by $1/\pi_y$ relative
to that earlier draft, so the normalized residual $\norm{v_k}/\varrho_k$ shrinks accordingly.
\end{remark}

\begin{definition}[Gradient generation units in the federated setting]
\label{def:ggu-fed}
In the federated instantiation of \cref{def:ggu} used throughout, a GGU is a \emph{worker}
$i\in[n_b]$ owning a shard $\mathcal{S}_i$ of size $n/n_b$. Before training the adversary fixes,
for every $i\in\mathcal{M}$, selection rates and a flipping kernel as in \cref{def:kernel}, hence
flip masses $u^i$; this produces once and for all a poisoned shard $\tilde{\mathcal{S}}_i$
(with $u^i=0$ and $\tilde{\mathcal{S}}_i=\mathcal{S}_i$ for honest units). At round $k$, unit $i$
draws a minibatch $B^i_k\subseteq\tilde{\mathcal{S}}_i$ and emits
\begin{equation}
g_i(\thk)=\frac{1}{|B^i_k|}\sum_{(x,\tilde y)\in B^i_k}\nabla_{\thk}\ell_c\big(f_{\thk}(x),\tilde y\big),
\label{eq:ggu-fed}
\end{equation}
$\tilde y$ denoting the (possibly altered) label carried by the sample. Equivalently, by
\cref{lem:message},
\begin{equation}
g_i(\thk)=\gradc(\thk)+\bar G(\thk)u^i+\xi_i,
\qquad \Ethk{\xi_i}=0,
\qquad \Ethk{\normsq{\xi_i}}=:\sigma_i^2(\thk),
\label{eq:ggu-fed-split}
\end{equation}
$\sigma_i^2(\thk)$ being the minibatch variance over $\tilde{\mathcal{S}}_i$.
\textbf{Convention.} $\Npk$ counts the \emph{corrupted units}, i.e.\ $\Npk:=|\mathcal{M}|=n_p$
for every $k$ --- not the number of altered labels drawn in the round, which is random and may
be zero for a corrupted unit. Under this convention the event $\mathcal{C}_k=\{\Npk\le f\}$ of
\cref{def:counts} is deterministic, holding for every $k$ as soon as $n_p\le f$.
\end{definition}

\begin{remark}[Two consequences]
\label{rem:ggu-fed}
\emph{(i)} A corrupted unit emits a \emph{partially} poisoned message at every round, rather than
an arbitrary vector at a random subset of rounds; the batch composition is fixed and only the
minibatch noise is random.
\emph{(ii)} The deviation of a corrupted message from the honest mean is proportional to the local
corruption rate, $\norm{\bar G(\thk)u^i}\le\varsigma_k\|u^i\|_1=\beta^i\varsigma_k$ with
$\varsigma_k$ as in \cref{eq:varsigma}: a small budget makes a corrupted unit \emph{permanently
stealthy} rather than rarely visible, which is what makes \cref{sec:stealth-selection} meaningful.
Part~0 and Part~A use only the first moment \cref{lem:message}; Part~B, whose rules select
\emph{realized} messages, uses \cref{eq:ggu-fed-split} and~(ii).
\end{remark}

\subsection{The three admissible sets}
\label{sec:three-sets}

Three distinct objects appear below; all are defined from the write access alone and \textbf{none presupposes any aggregation rule}.

\begin{definition}[Configurations, per-unit masses, aggregate masses]
\label{def:config}
\begin{align}
\mathfrak{U}_\beta&:=\Big\{(u^i)_{i=1}^{n_b}:u^i\ge0,\ \textstyle\sum_{z\neq y}u^i_{y,z}\le\pi_y\ \forall y,\ u^i=0\ \forall i\notin\mathcal{M},\ \tfrac1{n_b}\sum_i\|u^i\|_1\le\beta\Big\},
\label{eq:Ufrak}\\
\mathcal{U}^{\mathrm{loc}}&:=\Big\{u\ge0:\textstyle\sum_{z\neq y}u_{y,z}\le\pi_y\ \forall y,\ \|u\|_1\le\beta/\gamma\Big\},
\label{eq:Uloc}\\
\mathcal{U}_\beta&:=\Big\{u\ge0:\ \|u\|_1\le\beta,\ \textstyle\sum_{z\neq y}u_{y,z}\le\gamma\,\pi_y\ \forall y\Big\}.
\label{eq:Ubeta}
\end{align}
\end{definition}

\begin{remark}[What each set is for]
\label{rem:three-sets}
\begin{itemize}[noitemsep,leftmargin=*]
\item $\mathfrak{U}_\beta\subset(\R^{C(C-1)})^{n_b}$ is the set of \emph{whole configurations}: one flip-mass vector per unit, jointly constrained by the global budget. It is the adversary's genuine decision variable.
\item $\mathcal{U}^{\mathrm{loc}}\subset\R^{C(C-1)}$ is the set of \emph{single-unit} flip masses when the budget is split equally; it is nonempty because $\beta/\gamma\le1$ by \cref{lem:beta-bar}. It is the natural parameter when the aggregation rule sees individual messages (Part~B), and the one an implementation naturally exposes when the budget is stated per corrupted worker (\cref{rem:units}).
\item $\mathcal{U}_\beta\subset\R^{C(C-1)}$ is the set of \emph{aggregate} flip masses $\bar u:=\tfrac1{n_b}\sum_iu^i$. \cref{lem:hom-wlog} shows it is exactly the image of $\mathfrak{U}_\beta$ under averaging. \textbf{This is a statement about the budget, not about $\Agg$}: the aggregation rule enters only in \cref{prop:Ek-mean}, which identifies $\bar u$ with the induced bias when $\Agg$ is the mean.
\end{itemize}
\end{remark}

The three spaces are related as follows.

\begin{lemma}[Exact range of $\bar u$; homogeneity is without loss]
\label{lem:hom-wlog}
$\{\bar u:(u^i)\in\mathfrak{U}_\beta\}=\mathcal{U}_\beta$, and every $\bar u\in\mathcal{U}_\beta$ is attained by the homogeneous configuration $u^i=\bar u/\gamma$ for $i\in\mathcal{M}$, $u^i=0$ otherwise. In particular $\mathcal{U}_\beta=\gamma\,\mathcal{U}^{\mathrm{loc}}$.
\end{lemma}
\begin{proof}
($\subseteq$) $\|\bar u\|_1\le\tfrac1{n_b}\sum_i\|u^i\|_1\le\beta$ and $\sum_{z\neq y}\bar u_{y,z}=\tfrac1{n_b}\sum_{i\in\mathcal{M}}\sum_{z\neq y}u^i_{y,z}\le\gamma\pi_y$.
($\supseteq$) Put $u^i:=\bar u/\gamma$ on $\mathcal{M}$. Then $u^i\ge0$; capacity holds since $\sum_{z\neq y}u^i_{y,z}=\tfrac1\gamma\sum_{z\neq y}\bar u_{y,z}\le\pi_y$; the budget holds since $\tfrac1{n_b}\sum_i\|u^i\|_1=\|\bar u\|_1\le\beta$; and $\beta^i=\|u^i\|_1\le\beta/\gamma\le1$ by \cref{lem:beta-bar}. Comparing \cref{eq:Uloc,eq:Ubeta} coordinatewise gives $\mathcal{U}_\beta=\gamma\,\mathcal{U}^{\mathrm{loc}}$.
\end{proof}

\begin{remark}[Units: aggregate, local, and label counts]
\label{rem:units}
Three quantities are easily confused, and an implementation must keep them apart.
\begin{itemize}[noitemsep,leftmargin=*]
\item \emph{Aggregate.} $\bar u\in\mathcal{U}_\beta$ with $\|\bar u\|_1\le\beta$. This is the
variable of \cref{prop:Ek-mean}: $b_k=\bar G(\thk)\bar u$, with \textbf{no} extra factor.
\item \emph{Local.} $u^i\in\mathcal{U}^{\mathrm{loc}}$ with $\|u^i\|_1\le\beta/\gamma$. Under the
homogeneous deployment of \cref{lem:hom-wlog}, $\bar u=\gamma u^i$, hence
\begin{equation}
b_k=\gamma\,\bar G(\thk)\,u^i .
\label{eq:local-scope}
\end{equation}
A formulation whose decision variable is the per-worker $u^i$ --- the natural one when the
budget is stated as a fraction of a corrupted worker's own shard --- must therefore carry the
factor $\gamma$ explicitly. Omitting it compares a worker-level deviation with an
aggregate-level target and overstates the reachable set by $1/\gamma$.
\item \emph{Label counts.} The number of class-$y$ samples relabelled to $z$ across all corrupted
units is $n\,\bar u_{y,z}$ in aggregate units, equivalently $\gamma\,n\,u^i_{y,z}$ in local
units, and $\sum_{y,z}n\,\bar u_{y,z}=N_{\mathrm{flip}}\le\beta n$ recovers \cref{def:budget}.
\end{itemize}
\end{remark}

\begin{lemma}[Capacity redundancy]
\label{lem:Ubeta}
If $\beta\le\gamma\min_y\pi_y$ the capacity constraints in \cref{eq:Ubeta} are inactive and $\mathcal{U}_\beta=\beta\bar\Delta$, where $\bar\Delta:=\{u\ge0:\|u\|_1\le1\}=\mathrm{conv}(\{0\}\cup\{e_{y,z}\}_{z\neq y})$ is the \emph{corner simplex} and $e_{y,z}$ the canonical basis of $\R^{C(C-1)}$.
\end{lemma}
\begin{proof}
For admissible $u$ and any $y$, $\sum_{z\neq y}u_{y,z}\le\|u\|_1\le\beta\le\gamma\pi_y$. For the identification, $u\ge0$ with $\|u\|_1\le1$ writes $\sum_{z\neq y}u_{y,z}e_{y,z}+(1-\|u\|_1)\cdot0$, a convex combination of $\{0\}\cup\{e_{y,z}\}$; conversely every such combination is nonnegative with $\ell_1$ norm at most $1$.
\end{proof}

For balanced classes and small budgets the hypothesis holds: $C=10$, $n_b=10$, $n_p=3$,
$\gamma=0.3$ and $\pi_y=1/10$ give the threshold $\beta^{\mathrm{cap}}:=\gamma\min_y\pi_y=3\%$.
The vertex $0$ encodes unused budget, which is what makes $\bar\Delta$ a corner rather than a
standard simplex. Beyond that threshold the capacity constraints bind, $\mathcal{U}_\beta
\subsetneq\beta\bar\Delta$, and $\beta\mathcal{F}(\thk)$ is only an outer relaxation of the
reachable set; we write $s_\beta:=\beta/\beta^{\mathrm{cap}}$ for the \emph{saturation index} and
record in \cref{lem:knapsack,rem:saturated} exactly what this changes.

\begin{lemma}[Support function of $\mathcal{U}_\beta$: fractional knapsack]
\label{lem:knapsack}
For $c\in\R^{C(C-1)}$ set $\psi_y(c):=\max\{0,\max_{z\neq y}c_{y,z}\}$ and order the classes so
that $\psi_{(1)}\ge\cdots\ge\psi_{(C)}$. Then
\begin{equation}
\max_{u\in\mathcal{U}_\beta}\langle c,u\rangle=\sum_{y=1}^{C}\psi_{(y)}\,m^\star_{(y)},
\qquad
m^\star_{(y)}=\min\Big\{\gamma\pi_{(y)},\ \big(\beta-\textstyle\sum_{y'<y}m^\star_{(y')}\big)_+\Big\},
\label{eq:knapsack}
\end{equation}
i.e.\ the budget is poured greedily into the classes of highest gain, up to their capacities.
If $\beta\le\gamma\min_y\pi_y$ this reduces to $\beta\,\psi_{(1)}(c)$.
\end{lemma}
\begin{proof}
For $u\in\mathcal{U}_\beta$ put $m_y:=\sum_{z\neq y}u_{y,z}$; then $\langle c,u\rangle\le
\sum_ym_y\psi_y(c)$, and $m$ satisfies $0\le m_y\le\gamma\pi_y$, $\sum_ym_y\le\beta$. Conversely
any such $m$ is realized by $u_{y,z^\star_y}=m_y$ with $z^\star_y\in\argmax_{z\neq y}c_{y,z}$,
taking $m_y=0$ on the classes with $\psi_y(c)=0$. The relaxed problem
$\max\{\sum_ym_y\psi_y:0\le m_y\le\gamma\pi_y,\ \sum_ym_y\le\beta\}$ is a continuous knapsack: if a
feasible $m$ puts mass on a class $y$ while a class $y'$ with $\psi_{y'}>\psi_y$ is unsaturated,
transferring mass from $y$ to $y'$ is feasible and does not decrease the objective; iterating gives
\cref{eq:knapsack}. The last claim is the case where the first class absorbs the whole budget.
\end{proof}

\subsection{One-shot poisoning}
\label{sec:oneshot}

Based on this formalism, we can define the policy family which would be followed during the
label-poisoned training.

\begin{definition}[Per-round label-flipping family]
\label{def:qk-flip}
Writing $\ga^{(u)}(\thk)=(\ga^{i}(\thk))_{i\in\mathcal{M}}$ for the joint policy induced by a
configuration through \cref{def:kernel}, the \emph{per-round label-flipping family} is
\begin{equation}
\qk^{\mathrm{flip}}:=\big\{\ga^{(u)}(\thk):(u^i)\in\mathfrak{U}_\beta\big\}\ \subseteq\ \qk .
\label{eq:qk-flip}
\end{equation}
\end{definition}

\begin{definition}[Induced policy family; one-shot restriction]
\label{def:induced}
The adversary selects once, before training, a configuration $(u^i)\in\mathfrak{U}_\beta$ and a
trigger $\delta$, which fix the poisoned dataset. For each round $k$ the policy $\ga^i(\thk)$ is
then \emph{induced}: it is the law of the gradient of the fixed poisoned shard evaluated at
$\thk$. At the level of \emph{sequences} of policies,
\begin{equation}
\mathcal{Q}^{\mathrm{os}}:=\big\{\big(\ga^{(u)}(\thk)\big)_{k}:(u^i)\in\mathfrak{U}_\beta\big\}
\ \subsetneq\
\mathcal{Q}^{\mathrm{ad}}:=\prod_k\qk^{\mathrm{flip}} ,
\end{equation}
the one-shot family being the \emph{diagonal} of the adaptive one.
\end{definition}

$\qk$ permits a fresh policy at each round; one-shot poisoning does not. The restriction has
\textbf{no content at a fixed round} --- the projection of $\mathcal{Q}^{\mathrm{os}}$ onto round
$k$ is all of $\qk^{\mathrm{flip}}$ --- and bites only across rounds, since $(\ga^i(\thk))_k$ is
then a deterministic function of a single configuration; \cref{prop:oneshot-gap} records the
cost.

\begin{remark}[Parameterization of $\qk$ by the write access]
\label{rem:qk-param}
Label flipping restricts the admissible policies of \cref{def:ggu} to the image of the
configuration space, and the three sets of \cref{def:config} are exactly the three levels at
which \cref{eq:qk-flip} is indexed: \emph{(i)} $\mathfrak{U}_\beta$ parameterizes it in full;
\emph{(ii)} restricted to homogeneous configurations, $\mathcal{U}^{\mathrm{loc}}$ parameterizes
the policy of a \emph{single} corrupted unit, which is the level seen by an aggregation rule
(Part~B); \emph{(iii)} $\mathcal{U}_\beta$ parameterizes the induced bias under the mean,
$b_k=\bar G(\thk)\bar u$ (\cref{prop:Ek-mean}), and \cref{lem:hom-wlog} states that (ii) and
(iii) lose nothing. \cref{rem:units} gives the conversion factors between the three.
\end{remark}


\subsection*{Part A --- Mean aggregation}

The goal of this part is to derive an objective the attacker aims to optimize in order to find
the optimal trigger and attack policy, given their poisoning budget, target architecture and
learning setup (assuming the defender uses the Mean aggregation rule).

\subsection{Exact description of $\mathcal{E}_k$ under the mean}
\label{sec:Ek-mean}

\begin{proposition}[Exact description under the mean]
\label{prop:Ek-mean}
\textbf{Assume $\Agg$ is the mean.} Then, conditionally on $\thk$,
\begin{equation}
b_k:= \Ethk{h_k} - \gradc(\thk) = \bar G(\thk)\,\bar u\ \text{ with }\bar u\in\mathcal{U}_\beta,
\qquad
\mathcal{E}_k^{\mathrm{mean}}=\bar G(\thk)\,\mathcal{U}_\beta ,
\end{equation}
and if moreover $\beta\le\gamma\min_y\pi_y$, then $\mathcal{E}_k^{\mathrm{mean}}=\beta\,\mathcal{F}(\thk)$.
\end{proposition}

\begin{proof}
Linearity of the mean and \cref{lem:message} give $\Ethk{h_k}=\tfrac1{n_b}\sum_i\Ethk{g_i(\thk)}=\gradc(\thk)+\bar G(\thk)\bar u$, so $b_k=\bar G(\thk)\bar u$ by \cref{def:hk-decomposition}. Taking images as $(u^i)$ ranges over $\mathfrak{U}_\beta$ and using \cref{lem:hom-wlog} yields $\mathcal{E}_k^{\mathrm{mean}}=\bar G(\thk)\mathcal{U}_\beta$. Under the extra hypothesis \cref{lem:Ubeta} gives $\mathcal{U}_\beta=\beta\bar\Delta$, and since a linear map commutes with scaling and convex hulls, $\bar G(\beta\,\mathrm{conv}(\{0\}\cup\{e_{y,z}\}))=\beta\,\mathrm{conv}(\{0\}\cup\{\bar G_{y,z}\})=\beta\mathcal{F}(\thk)$.
\end{proof}

\begin{remark}[Hard flips]
\label{rem:hard-flips}
Every realization of a policy in \cref{def:kernel} is a hard relabeling; $\mathcal{F}(\thk)$ is the set of \emph{expected} gradients under randomized hard flips, around which the empirical gradient of a poisoned $n$-sample concentrates at rate $O(n^{-1/2})$. The convex relaxation is exact in expectation and is \emph{not} a soft-label relaxation.
\end{remark}

\begin{remark}[No robustness coefficient here]
\label{rem:mean-not-robust}
$\kappastar{f}(\mathrm{mean})=+\infty$: a single unit sending $g_1\to\infty$ drives the numerator of \cref{def:kappa_star} to infinity at fixed denominator. \cref{lem:bk-kappa} and \cref{prop:ball_inclusion} are therefore \emph{vacuous} for the mean and are used nowhere in Part~A. What bounds $\mathcal{E}_k^{\mathrm{mean}}$ is not the defense but solely the write access.
\end{remark}


\subsection{Geometry of $\mathcal{E}_k^{\mathrm{mean}}$}
\label{sec:geometry}

\begin{lemma}[Support function and radius]
\label{lem:support}
\textbf{Assume $\Agg$ is the mean}, and for $p\in\R^d$ set
$\psi_y(p):=\max\{0,\max_{z\neq y}\langle p,\bar G_{y,z}(\thk)\rangle\}$. Then
\begin{equation}
\max_{e\in\mathcal{E}_k^{\mathrm{mean}}}\langle p,e\rangle
=\sum_{y=1}^{C}\psi_{(y)}(p)\,m^\star_{(y)}\qquad\text{with the allocation \cref{eq:knapsack}},
\label{eq:support-general}
\end{equation}
and the exact radius $\varrho^{\mathrm{ex}}_k:=\max_{e\in\mathcal{E}_k^{\mathrm{mean}}}\norm{e}$
satisfies
\begin{equation}
\min\{\beta,\gamma\pi_{y^\star}\}\,\varsigma_k\ \le\ \varrho^{\mathrm{ex}}_k\ \le\ \beta\varsigma_k,
\qquad y^\star\in\argmax_y\max_{z\neq y}\norm{\bar G_{y,z}(\thk)} ,
\label{eq:radius-general}
\end{equation}
with $\varsigma_k$ as in \cref{eq:varsigma}. If $\beta\le\gamma\pi_{y^\star}$ --- in particular if
$\beta\le\gamma\min_y\pi_y$ --- the right-hand inequality is an equality, and under
$\beta\le\gamma\min_y\pi_y$ the support function reduces to
$\beta\max\{0,\max_{z\neq y}\langle p,\bar G_{y,z}(\thk)\rangle\}$.
\end{lemma}
\begin{proof}
$\mathcal{E}_k^{\mathrm{mean}}=\bar G(\thk)\mathcal{U}_\beta$ by \cref{prop:Ek-mean}, so
$\langle p,\bar G(\thk)u\rangle=\langle\bar G(\thk)^\top p,u\rangle$ and \cref{lem:knapsack}
applies with $c=\bar G(\thk)^\top p$, whose $\psi_y$ is the one above. For the radius,
$\norm{\bar G(\thk)u}\le\varsigma_k\|u\|_1\le\beta\varsigma_k$ gives the upper bound, while
$u=\min\{\beta,\gamma\pi_{y^\star}\}e_{y^\star,z^\star}$ is admissible and attains
$\min\{\beta,\gamma\pi_{y^\star}\}\varsigma_k$; the two coincide when $\beta\le\gamma\pi_{y^\star}$,
where \cref{lem:Ubeta} also gives the closed form.
\end{proof}

\begin{remark}[Saturated regime $s_\beta>1$: what changes and what does not]
\label{rem:saturated}
\emph{Unaffected.} Every statement using only that $\mathcal{U}_\beta$ is a compact convex polytope
containing $0$, and that $\mathcal{E}_k^{\mathrm{mean}}=\bar G(\thk)\mathcal{U}_\beta$: the first
identity of \cref{prop:Ek-mean}, \cref{lem:proj}, \cref{prop:dim}, \cref{rem:rank},
\cref{prop:ak-dist} (the QP simply keeps the two families of constraints of \cref{eq:Ubeta}),
\cref{sec:solution} in its entirety (\cref{eq:qp}, \cref{prop:danskin}), \cref{prop:oneshot-gap},
and all of Part~B, which uses only $\|u^i\|_1\le\beta/\gamma$ and $\varsigma_k$.

\emph{Affected, but only in the constants.} \emph{(i)} $\mathcal{E}_k^{\mathrm{mean}}\subsetneq
\beta\mathcal{F}(\thk)$: the corner-simplex description is an outer relaxation. \emph{(ii)} The
support function is the water-filling \cref{eq:support-general} instead of a single maximum, and the
radius obeys \cref{eq:radius-general}; $\beta\varsigma_k$ overestimates $\varrho^{\mathrm{ex}}_k$ by
at most $\max\{1,\beta/(\gamma\pi_{y^\star})\}\le s_\beta$. \emph{(iii)} (M1) and (M2) of
\cref{prop:certificates-mean} remain \emph{valid} certificates as stated, since
$\mathcal{E}_k^{\mathrm{mean}}\subseteq\beta\mathcal{F}(\thk)\subseteq\bar B(0,\beta\varsigma_k)$;
they are merely conservative, their sharp versions replacing $\beta\varsigma_k$ by
$\varrho^{\mathrm{ex}}_k$ and the maximum by \cref{eq:support-general}. (M3) is untouched.
\emph{(iv)} $\varrho_k=\beta\varsigma_k$ still satisfies (R1)--(R3) of \cref{sec:norm} and remains
the normalization we use; only its reading changes, from ``the largest reachable deviation'' to
``the largest deviation the budget could buy if it could be spent on a single class'', the gap being
at most $s_\beta$.

\emph{Genuinely lost.} \cref{prop:budget-match}: homogeneity in $\beta$ requires $\mathcal{U}_\beta$
to scale with $\beta$, which fails once capacities bind --- $\tfrac1\lambda\mathcal{U}_\beta$ then
depends on $\beta/\lambda$ \emph{and} on $\gamma\pi_y/\lambda$ separately, so $\lambda=\beta$ no
longer reduces exact feasibility to a budget-free condition.

\emph{Structural reading.} Since $\sum_{z\neq y}\bar u_{y,z}\le\gamma\pi_y$, the mass the aggregate
can carry \emph{out of a single source class} is capped at $\gamma\pi_{y_{\mathrm{source}}}$
irrespective of $\beta$. For a backdoor concentrated on one source class, budget beyond
$\gamma\pi_{y_{\mathrm{source}}}$ is therefore \emph{unusable}: $\mathcal{E}_k^{\mathrm{mean}}$
plateaus in the relevant directions, and the extra alterations either serve other classes or are
left unspent (the vertex $0$ of \cref{lem:Ubeta}). This is a prediction about the budget sweep, not
an artifact of the analysis.
\end{remark}

\begin{lemma}[Projection]
\label{lem:proj}
\textbf{Assume $\Agg$ is the mean.} Then \emph{(i)} $\mathcal{E}_k^{\mathrm{mean}}$ is nonempty, convex and compact; \emph{(ii)} for every $q\in\R^d$ the projection $P_k(q):=\argmin_{e\in\mathcal{E}_k^{\mathrm{mean}}}\norm{e-q}$ exists and is unique; \emph{(iii)} $P_k(q)$ is the unique $e^\star$ with $\langle q-e^\star,e-e^\star\rangle\le0$ for all $e\in\mathcal{E}_k^{\mathrm{mean}}$; \emph{(iv)} $q\mapsto\operatorname{dist}(q,\mathcal{E}_k^{\mathrm{mean}})^2$ is differentiable with gradient $2(q-P_k(q))$.
\end{lemma}
\begin{proof}
\emph{(i)} $0\in\mathcal{U}_\beta$ gives $0\in\mathcal{E}_k^{\mathrm{mean}}$; $\mathcal{U}_\beta$ is cut out by finitely many linear inequalities and bounded, hence a compact convex polytope, and linear images preserve convexity and compactness.
\emph{(ii)} Continuity on a compact set gives existence. If $e_1\neq e_2$ both attain the minimum value $r$, convexity gives $\tfrac{e_1+e_2}2\in\mathcal{E}_k^{\mathrm{mean}}$ and the parallelogram identity gives $\normsq{\tfrac{e_1+e_2}2-q}=\tfrac12\normsq{e_1-q}+\tfrac12\normsq{e_2-q}-\tfrac14\normsq{e_1-e_2}<r^2$, a contradiction.
\emph{(iii)} ($\Rightarrow$) For $t\in(0,1]$, $e^\star+t(e-e^\star)\in\mathcal{E}_k^{\mathrm{mean}}$, so $\normsq{q-e^\star}\le\normsq{q-e^\star}-2t\langle q-e^\star,e-e^\star\rangle+t^2\normsq{e-e^\star}$, whence $2\langle q-e^\star,e-e^\star\rangle\le t\normsq{e-e^\star}$; let $t\to0^+$. ($\Leftarrow$) $\normsq{q-e}=\normsq{q-e^\star}-2\langle q-e^\star,e-e^\star\rangle+\normsq{e-e^\star}\ge\normsq{q-e^\star}$; uniqueness by (ii).
\emph{(iv)} $P_k$ is $1$-Lipschitz: writing (iii) at $q_1$ with $e=P_kq_2$ and at $q_2$ with $e=P_kq_1$ and adding gives $\normsq{P_kq_1-P_kq_2}\le\langle q_1-q_2,P_kq_1-P_kq_2\rangle\le\norm{q_1-q_2}\norm{P_kq_1-P_kq_2}$. With $\Psi(q):=\normsq{q-P_k(q)}$, the bounds $\Psi(q+h)\le\normsq{q+h-P_kq}$ and $\Psi(q)\le\normsq{q-P_k(q+h)}$ give
$2\langle q-P_k(q+h),h\rangle+\normsq{h}\le\Psi(q+h)-\Psi(q)\le2\langle q-P_kq,h\rangle+\normsq{h}$,
and $\norm{P_k(q+h)-P_kq}\le\norm{h}$ makes both sides $2\langle q-P_kq,h\rangle+O(\normsq{h})$.
\end{proof}

\begin{proposition}[Dimension ceiling]
\label{prop:dim}
\textbf{Assume $\Agg$ is the mean.} $\mathcal{E}_k^{\mathrm{mean}}\subseteq\mathcal{G}_k:=\operatorname{Im}\bar G(\thk)$, a subspace of dimension at most $C(C-1)$. Hence $\mathcal{E}_k^{\mathrm{mean}}$ is Lebesgue-negligible in $\R^d$ whenever $C(C-1)<d$, and no ball description of it can be tight.
\end{proposition}
\begin{proof}
$\mathcal{E}_k^{\mathrm{mean}}=\bar G(\thk)\mathcal{U}_\beta\subseteq\operatorname{Im}\bar G(\thk)$ by \cref{prop:Ek-mean}, and $\operatorname{Im}\bar G(\thk)$ is spanned by the $C(C-1)$ columns. A proper subspace has empty interior and zero Lebesgue measure, whereas $\bar B(0,r)$ with $r>0$ has positive measure.
\end{proof}

\begin{remark}[Rank ceiling; to be checked experimentally]
\label{rem:rank}
Let $P^{\mathcal{G}}_k$ project orthogonally onto $\mathcal{G}_k$ and set
$\varpi_k:=\normsq{P^{\mathcal{G}}_kv_k}/\normsq{v_k}\in[0,1]$, the \emph{rank ratio}. Since
$\mathcal{E}_k^{\mathrm{mean}}\subseteq\mathcal{G}_k$, Pythagoras gives
\begin{equation}
a_k=\normsq{(I-P^{\mathcal{G}}_k)v_k}+\operatorname{dist}\big(P^{\mathcal{G}}_kv_k,\mathcal{E}_k^{\mathrm{mean}}\big)^2\ \ge\ (1-\varpi_k)\normsq{v_k}.
\label{eq:rank-ceiling}
\end{equation}
For $v_k$ uniform on the sphere $\mathbb{E}[\varpi_k]=\dim\mathcal{G}_k/d$, so the meaningful
comparison is $\varpi_k$ against $\dim\mathcal{G}_k/d$, not against $1$. Numerically, $\varpi_k$
should be evaluated with a truncated least-squares solve rather than a bare pseudo-inverse:
$\bar G(\thk)^\top\bar G(\thk)$ is typically ill-conditioned, and an untruncated $\mathrm{pinv}$
amplifies its smallest eigendirections enough to return values above $1$.
\end{remark}


\subsection{Feasibility under the mean}
\label{sec:feas-mean}

\begin{proposition}[$a$-feasibility as a projection]
\label{prop:ak-dist}
For any aggregation rule and every $k$,
$a_k=\inf_{\ga\in\qk^{\mathrm{flip}}}\normsq{b_k^\dagger}=\operatorname{dist}(v_k,\mathcal{E}_k)^2$,
the distance to a possibly nonconvex set. \textbf{If $\Agg$ is the mean}, the infimum is attained
at the unique $b_k^\star=P_k(v_k)$, computed by the convex QP
$\min_{u\in\mathcal{U}_\beta}\normsq{\bar G(\thk)u-v_k}$ in $O(C^2)$ variables.
\end{proposition}
\begin{proof}
$b_k^\dagger=b_k-v_k$ by \cref{prop:biais-relationship} and $b_k$ ranges exactly over
$\mathcal{E}_k$ as $\ga$ ranges over $\qk^{\mathrm{flip}}$, by
\cref{def:feasible-attack-set,def:qk-flip}; take the infimum. Under the mean,
\cref{prop:Ek-mean} identifies $\mathcal{E}_k$ and \cref{lem:proj} gives existence, uniqueness
and the QP form.
\end{proof}

\begin{proposition}[Certificates of infeasibility under the mean]
\label{prop:certificates-mean}
\textbf{Assume $\Agg$ is the mean}, and let $v_k\neq0$. Each of the following implies $a_k>0$:
\begin{align}
\textbf{(M1)}\quad&\norm{v_k}>\beta\varsigma_k
&&\text{(radius, \cref{lem:support})},\\
\textbf{(M2)}\quad&\exists\,p\in\R^d:\ \langle p,v_k\rangle>\beta\max\{0,\max_{z\neq y}\langle p,\bar G_{y,z}(\thk)\rangle\}
&&\text{(separating hyperplane)},\\
\textbf{(M3)}\quad&\varpi_k<1
&&\text{(rank, \cref{rem:rank})}.
\end{align}
The sharp forms of (M1)--(M2) replace $\beta\varsigma_k$ by $\varrho^{\mathrm{ex}}_k$ and the
right-hand side of (M2) by \cref{eq:support-general}; the two coincide under
$\beta\le\gamma\min_y\pi_y$.
\end{proposition}
\begin{proof}
By \cref{prop:Ek-mean} and $\mathcal{U}_\beta\subseteq\beta\bar\Delta$ we have
$\mathcal{E}_k^{\mathrm{mean}}\subseteq\beta\mathcal{F}(\thk)\subseteq\bar B(0,\beta\varsigma_k)$,
the last inclusion because every element of $\beta\mathcal{F}(\thk)$ is a convex combination of
$0$ and the $\beta\bar G_{y,z}$, all of norm at most $\beta\varsigma_k$. Hence
$v_k\in\mathcal{E}_k^{\mathrm{mean}}$ implies $\norm{v_k}\le\beta\varsigma_k$ and, for every $p$,
$\langle p,v_k\rangle\le\max_{e\in\beta\mathcal{F}(\thk)}\langle p,e\rangle
=\beta\max\{0,\max_{z\neq y}\langle p,\bar G_{y,z}\rangle\}$; contrapose for (M1)--(M2). (M3) is
\cref{eq:rank-ceiling}. The sharp forms follow from \cref{lem:support} applied to
$\mathcal{E}_k^{\mathrm{mean}}$ itself.
\end{proof}

\begin{remark}[Exact feasibility is generically unattainable]
\label{rem:generic-infeasible}
(M3) is severe. By \cref{prop:dim}, $\dim\mathcal{G}_k\le C(C-1)\lll d$, so $\varpi_k<1$ for every $v_k$ outside a subspace of measure zero: \textbf{$a_k>0$ generically}, and exact feasibility ($v_k\in\mathcal{E}_k$) is not an attainable target but a degenerate limiting case. The object of $(\mathrm{P}^{\mathrm{mean}})$ is accordingly to \emph{minimize} $a_k$, not to annul it, and the normalization $\varrho_k$ of \cref{sec:norm} is what makes the achieved value interpretable. Statements of the form ``exact feasibility iff $v_k\in\mathcal{E}_k$'' are to be read as characterizations, not as goals.
\end{remark}

None of (M1)--(M3) involves $\Agg$, consistently with \cref{rem:mean-not-robust}. In particular proximity of $\gradp(\thk)$ to $\gradc(\thk)$ is \emph{not} the relevant criterion, since $\gradc(\thk)$ is the single vertex $0$ of $\mathcal{E}_k^{\mathrm{mean}}$.


\subsection{The trigger enters only through $v_k$}
\label{sec:trigger-vk}

Let $T_\delta(x)=x+\delta$, $\norm{\delta}_\infty\le\varepsilon$. By \cref{eq:bd-Dp,eq:bd-Lp-split}, with $\lambda$ the mixture weight,
\begin{equation}
v_k(\delta)=\gradp[\delta](\thk)-\gradc(\thk)
=\lambda\Big(\nabla_\theta\mathbb{E}_{X\sim\mathcal D_s}\big[\ell_c\big(f_{\thk}(T_\delta(X)),y_{\mathrm{target}}\big)\big]-\gradc(\thk)\Big).
\label{eq:vk-delta}
\end{equation}
The two decision variables act on disjoint objects: the configuration $(u^i)$ parameterizes $\mathcal{E}_k$, the trigger $\delta$ parameterizes $v_k$. This is the structural reason the reachable set is trigger-free.

\begin{remark}[$v_k$ does not vanish at $\delta=0$]
\label{rem:no-collapse}
At $\delta=0$, \cref{eq:vk-delta} still compares the gradient of assigning $y_{\mathrm{target}}$ to clean source-class inputs with $\gradc(\thk)$: it asks the model to classify the same inputs under two labels, so $\norm{v_k}$ is large. Minimizing $a_k$ therefore does \emph{not} collapse towards the null trigger; the efficacy term of \cref{eq:P} serves the different purpose of making the backdoor learnable.
\end{remark}

\begin{proposition}[Budget matching]
\label{prop:budget-match}
\textbf{Assume $\Agg$ is the mean and $\beta\le\gamma\min_y\pi_y$.} For all $\lambda,\beta>0$,
$\operatorname{dist}(v_k,\beta\mathcal{F}(\thk))=\lambda\operatorname{dist}\big(\tfrac1\lambda v_k,\tfrac\beta\lambda\mathcal{F}(\thk)\big)$,
so the geometry depends on $(\lambda,\beta)$ only through $\beta/\lambda$, and $\lambda=\beta$ reduces exact feasibility to the budget-free condition $\tfrac1\lambda v_k\in\mathcal{F}(\thk)$.
\end{proposition}
\begin{proof}
Homogeneity of the distance, $\operatorname{dist}(\lambda q,\lambda S)=\lambda\operatorname{dist}(q,S)$, with $q=\tfrac1\lambda v_k$ and $S=\tfrac\beta\lambda\mathcal{F}(\thk)$.
\end{proof}

Since $\tfrac1\lambda v_k$ is $\lambda$-free by \cref{eq:vk-delta}, the expert must be trained at
the poisoning rate the write access affords: if $\lambda\neq\beta$ the residual retains a
component that no configuration cancels and no trigger removes. The hypothesis is essential: in
the saturated regime the reduction fails, as recorded in \cref{rem:saturated}.


\subsection{Normalization}
\label{sec:norm}

$\norm{v_k}$ and $\varsigma_k$ vary by orders of magnitude along the trajectory, so $\mathbb{E}_k[a_k]$ would be dominated by the rounds carrying the largest gradients. A per-round scale $s_k>0$ is needed; we require it to (\textbf{R1}) equalize summands across $k$, (\textbf{R2}) not depend on $\delta$ \emph{given the reference trajectory} of \cref{sec:ref-traj}, so as not to distort the landscape being optimized, and (\textbf{R3}) be computable from quantities already assembled.

\begin{proposition}[Saturation of the scale-invariant criterion]
\label{prop:saturation}
Let $A\subset\R^d$ be convex compact with $0\in A$, $e\in\R^d$ with $\norm{e}=1$, and $r(t):=\operatorname{dist}(te,A)^2/t^2$. Then $r$ is nondecreasing on $(0,\infty)$, vanishes on $(0,t^\star]$ with $t^\star:=\max\{t:te\in A\}$, and $r(t)\to1$ with $r'(t)=O(t^{-2})$.
\end{proposition}
\begin{proof}
For $s<t$ and $e_t:=P_A(te)$, convexity and $0\in A$ give $\tfrac ste_t\in A$, so $\operatorname{dist}(se,A)\le\norm{se-\tfrac ste_t}=\tfrac st\operatorname{dist}(te,A)$; squaring gives monotonicity. For large $t$, $\operatorname{dist}(te,A)=t-h_A(e)+O(1/t)$ with $h_A$ the support function of \cref{lem:support}, so $r(t)=1-2h_A(e)/t+O(t^{-2})$.
\end{proof}

Both candidate scales leave the criterion equal to zero throughout $\mathcal{E}_k$ --- that is a
property of $a_k$, not of the normalization. What distinguishes them is the behaviour
\emph{outside}: taking $s_k=\norm{v_k}$ violates (R2) and, by \cref{prop:saturation}, saturates
at $1$ with gradient $O(t^{-2})$, so a trigger far outside $\mathcal{E}_k$ receives no usable
signal to return. We use instead the budget scale
\begin{equation}
\varrho_k:=\beta\varsigma_k=\beta\max_{z\neq y}\norm{\bar G_{y,z}(\thk)},
\label{eq:rho}
\end{equation}
which measures the residual in units of what the budget can buy, is $\delta$-free once the
trajectory is frozen, and grows quadratically outside $\mathcal{E}_k$ instead of saturating. By
\cref{eq:radius-general}, $\varrho_k$ \emph{upper bounds} the exact radius
$\varrho^{\mathrm{ex}}_k$, with equality iff $\beta\le\gamma\pi_{y^\star}$; a value below $1$
therefore means the residual is smaller than the largest deviation the budget could buy if it
could be spent on a single class, and in the unsaturated regime, smaller than the largest
reachable deviation itself.


\subsection{The attacker's problem under the mean}
\label{sec:attacker-problem}

\begin{equation}
\begin{aligned}
(\mathrm{P}^{\mathrm{mean}})\quad
\min_{\norm{\delta}_\infty\le\varepsilon,\ (u^i)\in\mathfrak{U}_\beta}\
&\mathbb{E}_{k\sim\tau([k_0,K])}\!\left[\frac{\normsq{\bar G(\bar\theta_k)\bar u-v_k(\delta)}}{\varrho_k^{2}}\right]\\
&+\kappa\,\mathbb{E}_{k}\Big[\mathbb{E}_{X\sim\mathcal{D}_s}\big[\ell_c\big(f_{\bar\theta_k}(T_\delta(X)),y_{\mathrm{target}}\big)\big]\Big]
\end{aligned}
\label{eq:P}
\end{equation}
\begin{equation*}
\text{s.t.}\quad
\bar\theta_{k+1}=\bar\theta_k-\ek\,\gradp[\delta](\bar\theta_k),
\qquad
\lambda=\beta ,
\end{equation*}
where $\kappa$ calibrates stealth against strength. No acceptance constraint appears, by
\cref{rem:mean-not-robust}. The decision variable is the pair $(\delta,\bar u)$ jointly; if an
implementation parameterizes by the per-worker $u^i$ instead, the first term reads
$\normsq{\gamma\bar G(\bar\theta_k)u^i-v_k(\delta)}/\varrho_k^2$ by \cref{eq:local-scope}.

\begin{proposition}[Per-round feasibility is a relaxation]
\label{prop:oneshot-gap}
Let $J(u)$ be the first term of \cref{eq:P}. Then $\inf_{(u^i)\in\mathfrak{U}_\beta}J(u)\ge\mathbb{E}_k[a_k/\varrho_k^2]$, with equality iff some configuration attains the per-round infimum for almost every $k$.
\end{proposition}
\begin{proof}
For fixed $(u^i)$ the integrand dominates its pointwise infimum, which is $a_k/\varrho_k^2$ by \cref{prop:ak-dist}; take expectations, then the infimum on the left.
\end{proof}

Thus $a_k$, defined round by round, is a \emph{lower bound} on what a one-shot adversary
achieves; the gap is exactly the price of $\mathcal{Q}^{\mathrm{os}}$ against
$\mathcal{Q}^{\mathrm{ad}}$ in \cref{def:induced}, i.e.\ of committing the dataset once.

\begin{proposition}[Structure]
\label{prop:structure}
\emph{(i)} $(\mathrm{P}^{\mathrm{mean}})$ is a discrete-time optimal control problem with \emph{static} control $(\delta,\bar u)$; the only nesting is the trajectory recursion and there is no inner $\argmin$ over $\theta$.
\emph{(ii)} $\mathcal{U}_\beta$ and $\{\norm{\delta}_\infty\le\varepsilon\}$ are compact convex; the objective is continuous in $(\delta,\bar u)$ whenever $\delta\mapsto\gradp[\delta]$ and $\delta\mapsto\mathcal{L}_{bd}^\delta$ are, so the minimum is attained.
\emph{(iii)} Writing $V(\delta):=\min_uJ(\delta,u)$ is \emph{not} a decoupling: $\min_\delta\min_u=\min_{(\delta,u)}$ is an identity; optimizing $\delta$ against any surrogate other than $\min_uJ$ would be.
\emph{(iv)} By \cref{lem:hom-wlog} the configuration may be taken homogeneous, and one-shot poisoning makes $\bar u$ shared across rounds, so the inner problem is a \emph{single} QP coupling the whole trajectory.
\end{proposition}

\begin{remark}[Two admissible solvers, and what distinguishes them]
\label{rem:solver}
By \cref{prop:structure}(iii) the problem is joint in $(\delta,\bar u)$, and two solvers are
consistent with it.
\emph{(a) Exact inner solve.} At each outer step, solve the QP of \cref{eq:qp} for
$\bar u^\star(\delta)$, detach it, and take one projected-gradient step on $\delta$ using
\cref{prop:danskin}. The reported objective is then exactly $\mathbb{E}_k[a_k/\varrho_k^2]$ at
the current $\delta$, and the hypergradient is the true one.
\emph{(b) Simultaneous descent.} Step $\delta$ and $\bar u$ together, projecting $\bar u$ onto
$\mathcal{U}_\beta$ after each step. This is cheaper and remains a descent method on the joint
objective, but $\bar u$ lags $\delta$: the reported value \emph{overstates}
$\mathbb{E}_k[a_k/\varrho_k^2]$, and \cref{prop:danskin} does not apply, so the $\delta$-gradient
is not the hypergradient of $V$.
The two are reconciled by reporting, alongside the running objective, the value attained by the
exact inner solve on the same $v_k$; the gap measures the lag and is the quantity to monitor.
\end{remark}


\mbcomment{Since exact feasibility notion is a quite strong notion, it could be interesting to obtain a similar optimisation problem starting from the feasibility in alignment notion.}

\subsection{Reference trajectory}
\label{sec:ref-traj}

$\bar\theta(\delta)$ in \cref{eq:P} is the expected SGD trajectory on $\Lp$, i.e.\ the expert trajectory. Recall from \cref{def:hk-decomposition} that $h_k=\gradp(\thk)+b_k^\dagger+\xi_k$ with $\Ethk{\xi_k}=0$.

\begin{proposition}[Self-consistency]
\label{prop:selfcons}
Assume $\gradp$ is $L$-Lipschitz, $\ek\le\eta$, $\sup_k\norm{b_k^\dagger}\le\epsilon_b$ and $\bar\sigma:=\sup_k\mathbb{E}\norm{\xi_k}$. Then
\begin{equation}
\mathbb{E}\,\norm{\thk-\bar\theta_k(\delta)}\ \le\ \frac{\epsilon_b+\bar\sigma}{L}\Big[(1+\eta L)^k-1\Big].
\label{eq:selfcons}
\end{equation}
In particular, if the aggregation is noiseless ($\xi_k\equiv0$) and the attack exact ($b_k^\dagger\equiv0$), then $\thk=\bar\theta_k(\delta)$ for all $k$.
\end{proposition}
\begin{proof}
Let $e_k:=\thk-\bar\theta_k$. Subtracting the two recursions,
$e_{k+1}=e_k-\ek\big(\gradp(\thk)-\gradp(\bar\theta_k)\big)-\ek\big(b_k^\dagger+\xi_k\big)$,
so pathwise $\norm{e_{k+1}}\le(1+\eta L)\norm{e_k}+\eta\big(\norm{b_k^\dagger}+\norm{\xi_k}\big)$. Taking expectations and iterating from $e_0=0$ gives $\mathbb{E}\norm{e_k}\le\eta(\epsilon_b+\bar\sigma)\sum_{l<k}(1+\eta L)^l$, which is \cref{eq:selfcons}. The last claim is the case $\epsilon_b=\bar\sigma=0$.
\end{proof}

\begin{remark}[What self-consistency does and does not say]
\label{rem:selfcons-scope}
\cref{eq:selfcons} bounds $\mathbb{E}\norm{\thk-\bar\theta_k}$, not $\norm{\Ethk{\thk}-\bar\theta_k}$: since $\gradp$ is not affine, $\mathbb{E}[\gradp(\thk)]\neq\gradp(\mathbb{E}[\thk])$ and no statement about the mean iterate follows. Moreover the aggregation noise $\bar\sigma$ \emph{does not vanish} when the attack is exact, so the reference trajectory is the fixed point of $(\mathrm{P}^{\mathrm{mean}})$ only at the level of the expected update. The bound is exponential in $k$ and is a consistency statement rather than a numerical guarantee; the direct measurement of item~5 in \cref{sec:todo} is what certifies it in practice.
\end{remark}

Two approximations must be distinguished. \textbf{(T1)} evaluating on $\bar\theta(\delta)$ rather than on the victim's iterates, controlled by \cref{eq:selfcons}. \textbf{(T2)} not re-recording $\bar\theta$ as $\delta$ varies: the exact hypergradient obeys $p_k=\nabla_\theta\ell_k+(I-\ek\nabla^2\Lp(\bar\theta_k))^{\!\top}p_{k+1}$ with $p_{K+1}=0$, and (T2) sets $p\equiv0$; if $\norm{I-\ek\nabla^2\Lp(\bar\theta_k)}\le\rho<1$ and $\sup_k\norm{\nabla_\theta\ell_k}\le G_\ell$, the error is at most $K\eta\bar BG_\ell/(1-\rho)$ with $\bar B:=\sup_k\norm{\partial_\delta\gradp[\delta](\bar\theta_k)}$. (T2) is removed by one outer loop: re-record $\bar\theta(\delta)$ and re-solve.

\begin{remark}[Momentum]
\label{rem:momentum}
\cref{prop:selfcons} is stated for the plain SGD update \cref{eq:sgd}. With momentum the state becomes $(\theta,m)$ and both the recursion and the adjoint of (T2) must be written on that pair; the argument is unchanged in structure.
\end{remark}


\subsection{Solution}
\label{sec:solution}

By \cref{lem:hom-wlog} the configuration may be taken homogeneous, so the decision reduces to $\bar u\in\mathcal{U}_\beta$, shared across rounds:
\begin{equation}
u^{\star}(\delta)=\argmin_{u\in\mathcal{U}_\beta}\ u^{\top}Q\,u-2(c^{\delta})^{\top}u,
\qquad
Q=\sum_k\omega_k\bar G_k^{\top}\bar G_k,
\quad
c^{\delta}=\sum_k\omega_k\bar G_k^{\top}v_k(\delta),
\label{eq:qp}
\end{equation}
with $\bar G_k:=\bar G(\bar\theta_k)$ and $\omega_k:=1/(\varrho_k^2|[k_0,K]|)$. Under (T2) the matrix $Q$ is $\delta$-free and assembled once; only $c^\delta$ changes along the optimization. The dimension is $O(C^2)$, independent of $K$ and of $d$.

\begin{proposition}[Hypergradient]
\label{prop:danskin}
\textbf{Assume $\Agg$ is the mean and (T2)}, so $\mathcal{U}_\beta$ is compact convex and $\delta$-independent and $\bar G_k$ is $\delta$-independent. If $u^\star(\delta)$ is unique, $V(\delta):=\min_{u\in\mathcal{U}_\beta}J(\delta,u)$ is differentiable at $\delta$ with $\nabla V(\delta)=\nabla_\delta J(\delta,u^\star(\delta))$, $u^\star$ \emph{detached}.
\end{proposition}
\begin{proof}
Danskin's theorem: the feasible set is compact and $\delta$-independent and $J$ is $C^1$ in $\delta$ with continuous partial gradient; uniqueness of the minimizer makes the directional derivative linear in the direction.
\end{proof}

Uniqueness holds when $Q\succ0$; otherwise $+\epsilon\normsq{u}$ restores it at $O(\epsilon)$ bias and makes $V$ of class $C^1$. Since $\bar G_k$ is computed once per checkpoint and reused, each outer step costs one backward pass for $v_k(\delta)$ per round plus one small QP. Projected gradient descent on $\delta$ with step $\le1/L$ converges to a projected stationary point at rate $O(1/\sqrt T)$.